% C015 椭圆焦点三角形四心坐标性质——多角度证明
% 约定：P=(x_0,y_0), y_0\ne0, O=(0,0), F_1=(-c,0), F_2=(c,0)。
\begin{proofbox}
\textbf{准备：只用一次椭圆定义，推出两个焦半径公式}

记 $d_1=PF_1$、$d_2=PF_2$。椭圆第一定义给出
\[
 d_1+d_2=2a.
\]
把两个距离\textbf{平方相减}，不用解根式：
\begin{align*}
 d_1^2-d_2^2
 &=\bigl[(x_0+c)^2+y_0^2\bigr]
   -\bigl[(x_0-c)^2+y_0^2\bigr]\\
 &=4cx_0.
\end{align*}
利用平方差公式：
\[
 (d_1-d_2)(d_1+d_2)=4cx_0
 \quad\Longrightarrow\quad
 d_1-d_2=\frac{2cx_0}{a}=2ex_0.
\]
由和、差得到
\[
 \boxed{PF_1=a+ex_0},\qquad
 \boxed{PF_2=a-ex_0}.
\]
又因为底边 $F_1F_2=2c$，三角形的半周长与面积为
\[
 \boxed{s=a+c},\qquad
 \boxed{S=\frac12(2c)|y_0|=c|y_0|}.
\]
\textit{理解：后续所有简化，都来源于两个对称焦点和一个恒定的焦半径之和。}

\tcbline
\textbf{一、重心 $G$——两种方法}

\textbf{方法 G1（最快：三点坐标平均）。}
\[
 G=\frac{P+F_1+F_2}{3}
 =\left(\frac{x_0-c+c}{3},\frac{y_0}{3}\right)
 =\boxed{\left(\frac{x_0}{3},\frac{y_0}{3}\right)}.
\]

\textbf{方法 G2（几何中线：一眼看出比例）。}
底边中点恰好是 $O=(0,0)$，因此 $PO$ 是中线。
重心把中线按\textbf{从顶点起 $2:1$} 分割，即 $PG:GO=2:1$。
所以 $\overrightarrow{OG}=\tfrac13\overrightarrow{OP}$，
仍得 $G=(x_0/3,y_0/3)$。

\tcbline
\textbf{二、内心 $I$——三种方法；第一种最值得记忆}

\textbf{方法 I1（最快、最直觉：切线长相等 $+$ 面积）。}

设内切圆与底边 $F_1F_2$ 相切于 $T$，
与 $PF_1,PF_2$ 分别相切于 $U,V$。
由同一点向圆所引的两条切线等长，
\[
 F_1T=F_1U,\qquad F_2T=F_2V,\qquad PU=PV.
\]
故
\begin{align*}
 PF_1-PF_2
 &=(PU+UF_1)-(PV+VF_2)\\
 &=F_1T-F_2T.
\end{align*}
因 $F_1=(-c,0),F_2=(c,0)$，且 $T$ 在底边内部，
\[
 F_1T=c+x_T,\qquad F_2T=c-x_T,
\]
两者之差就是 $2x_T$。内心在切点的垂线上，故 $x_I=x_T$，
\[
 \boxed{x_I=\frac{PF_1-PF_2}{2}=ex_0}.
\]
接着用面积 $S=sr$：
\[
 c|y_0|=(a+c)r
 \quad\Longrightarrow\quad
 r=\frac{c|y_0|}{a+c}.
\]
内心和 $P$ 在底边同侧，所以 $y_I$ 与 $y_0$ 同号，从而
\[
 \boxed{y_I=\frac{cy_0}{a+c}}.
\]
\textbf{记忆口令：底边切点定横坐标；面积除以半周长定纵坐标。}

\medskip
\textbf{方法 I2（代数通法：对边长度加权平均）。}

一般三角形的内心坐标等于\textbf{各顶点坐标乘其对边长，再除以周长}。
在本题中，顶点 $P,F_1,F_2$ 的对边依次为
$2c,d_2,d_1$，千万不要颠倒：
\[
 I=\frac{2cP+d_2F_1+d_1F_2}{2c+d_1+d_2}.
\]
分母为 $2(a+c)$，横坐标为
\begin{align*}
 x_I
 &=\frac{2cx_0-c(a-ex_0)+c(a+ex_0)}{2(a+c)}\\
 &=\frac{2cx_0(1+e)}{2a(1+e)}
 =\boxed{ex_0}.
\end{align*}
纵坐标为
\[
 y_I=\frac{2cy_0}{2(a+c)}
 =\boxed{\frac{cy_0}{a+c}}.
\]
\textit{这是可推广到任意三角形的代数方法；上面的方法 I1 则更容易现场还原。}

\medskip
\textbf{方法 I3（角平分线定理 $+$ 直线分点）。}

设顶角 $\angle F_1PF_2$ 的内角平分线与底边相交于 $D=(d,0)$。
由角平分线定理，
\[
 \frac{F_1D}{DF_2}=\frac{PF_1}{PF_2}
 =\frac{a+ex_0}{a-ex_0}.
\]
因为 $F_1D=c+d$、$DF_2=c-d$，交叉相乘得到
\[
 d=c\frac{PF_1-PF_2}{PF_1+PF_2}
 =c\frac{2ex_0}{2a}=e^2x_0.
\]
于是 $D=(e^2x_0,0)$。内心 $I$ 位于 $PD$ 上；
由面积法已经知道 $y_I/y_0=c/(a+c)$，故按直线上的同比例分点，
\begin{align*}
 x_I
 &=d+\frac{y_I}{y_0}(x_0-d)\\
 &=e^2x_0+\frac{c}{a+c}(1-e^2)x_0\\
 &=\left[e^2+\frac{e}{1+e}(1-e^2)\right]x_0
 =\boxed{ex_0}.
\end{align*}
\textit{额外收获：顶角平分线落到底边的点 $D$ 的横坐标是 $e^2x_0$；
而内心横坐标是 $ex_0$。两者不要混淆。}

\tcbline
\textbf{三、垂心 $H$——两种直接证明和一种联动证明}

\textbf{方法 H1（最快直接法：一条竖直高线 $+$ 向量垂直）。}

底边 $F_1F_2$ 在 $x$ 轴上，故过 $P$ 的高线为
\[
 \boxed{x=x_0}.
\]
设 $H=(x_0,h)$。过 $F_1$ 的高线与 $PF_2$ 垂直，
也就是向量 $\overrightarrow{F_1H}=(x_0+c,h)$
与 $\overrightarrow{F_2P}=(x_0-c,y_0)$ 点积为零：
\[
 (x_0+c)(x_0-c)+hy_0=0.
\]
所以
\[
 \boxed{h=\frac{c^2-x_0^2}{y_0}},\qquad
 \boxed{H=\left(x_0,\frac{c^2-x_0^2}{y_0}\right)}.
\]
\textit{这个写法比使用斜率更稳健：当 $x_0=c$ 时，$PF_2$ 竖直、斜率不存在，
点积法仍能直接使用。}

\medskip
\textbf{方法 H2（解析几何：垂直斜率法）。}

$PF_2$ 的方向向量是 $(x_0-c,y_0)$，其中 $y_0\ne0$。
与之垂直且经过 $F_1$ 的高线可写为
\[
 y=-\frac{x_0-c}{y_0}(x+c).
\]
令 $x=x_0$ 得
\[
 y_H=-\frac{(x_0-c)(x_0+c)}{y_0}
 =\frac{c^2-x_0^2}{y_0}.
\]
注意：上式采用的是\textbf{高线的斜率}，不会涉及不存在的 $PF_2$ 斜率。

\medskip
\textbf{方法 H3（最快联动：欧拉线）。}

一旦外心 $O'=(0,t)$ 已经求得，重心 $G=P/3$ 已知，
由一般三角形的欧拉线公式
\[
 H=3G-2O'=P-2O'
\]
立刻得到
\begin{align*}
 H
 &=\left(x_0,y_0-\frac{x_0^2+y_0^2-c^2}{y_0}\right)\\
 &=\boxed{\left(x_0,\frac{c^2-x_0^2}{y_0}\right)}.
\end{align*}
这个公式还可以不依赖记忆直接核验：令 $H=P-2O'$，
显然 $H$ 与 $P$ 同横坐标；并且
\[
 \overrightarrow{F_1H}\cdot\overrightarrow{F_2P}
 =x_0^2-c^2+y_0^2-2ty_0=0,
\]
最后一步正是外心方程 $2ty_0=x_0^2+y_0^2-c^2$。
两条高线都满足，所以 $H$ 确为垂心。

\tcbline
\textbf{四、外心 $O'$——两种方法}

\textbf{方法 O1（最快：底边中垂线 $+$ 两段距离相等）。}

$F_1F_2$ 的中点是 $O=(0,0)$，底边是水平线，所以其垂直平分线为 $x=0$。
因此\textbf{对任意非退化的焦点三角形}，外心必为 $O'=(0,t)$。
由 $O'F_2=O'P$，
\[
 c^2+t^2=x_0^2+(y_0-t)^2.
\]
展开、约去 $t^2$：
\[
 2y_0t=x_0^2+y_0^2-c^2.
\]
所以
\[
 \boxed{O'=\left(0,\frac{x_0^2+y_0^2-c^2}{2y_0}\right)}.
\]
\textbf{重要：外心在 $y$ 轴上根本不要求 $PF_1=PF_2$。}

\medskip
\textbf{方法 O2（圆方程：经过两个焦点的圆）。}

把三角形的外接圆写成一般式
\[
 x^2+y^2+Dx+Ey+F=0.
\]
代入 $F_1(-c,0)$ 与 $F_2(c,0)$，分别有
\[
 c^2-Dc+F=0,\qquad c^2+Dc+F=0.
\]
两式相减得 $D=0$，再得 $F=-c^2$。
将 $P(x_0,y_0)$ 代入，得到
\[
 E=\frac{c^2-x_0^2-y_0^2}{y_0}.
\]
该圆圆心是 $(-D/2,-E/2)$，恰好为
\[
 O'=\left(0,\frac{x_0^2+y_0^2-c^2}{2y_0}\right).
\]
\textit{顺便得到外接圆的简洁方程：}
\[
 \boxed{x^2+y^2-2y_{O'}y-c^2=0}.
\]

\tcbline
\textbf{五、把椭圆方程再用一次：检验与特殊位置}

从
\[
 \frac{x_0^2}{a^2}+\frac{y_0^2}{b^2}=1
 \quad\Longrightarrow\quad
 x_0^2=a^2-\frac{a^2}{b^2}y_0^2
\]
及 $a^2=b^2+c^2$ 可得
\[
 \boxed{y_{O'}=\frac{b^2}{2y_0}-\frac{c^2y_0}{2b^2}},
 \qquad
 \boxed{y_H=\frac{a^2y_0}{b^2}-\frac{b^2}{y_0}}.
\]
这两式不必作为主要记忆对象，但可以用于核对符号或研究极限。

\textbf{位置判断 1：} $x_0=\pm c$ 时 $H=F_2$ 或 $H=F_1$，
因为相应的焦点处有直角。

\textbf{位置判断 2：} $x_0^2+y_0^2=c^2$ 时 $O'=O$，
此时 $P$ 在以 $F_1F_2$ 为直径的圆上，$\angle F_1PF_2=90^\circ$。

\textbf{位置判断 3：} 椭圆顶点 $P=(0,\pm b)$ 时，三角形是等腰三角形，
四心都在 $y$ 轴上：
\[
 G=\left(0,\frac{y_0}{3}\right),\quad
 I=\left(0,\frac{cy_0}{a+c}\right),\quad
 H=\left(0,\frac{c^2}{y_0}\right),\quad
 O'=\left(0,\frac{b^2-c^2}{2y_0}\right).
\]
四心不一定重合；只有焦点三角形恰为等边三角形时才会重合。
\end{proofbox}
